% Part: normal-modal-logic
% Chapter: syntax-and-semantics
% Section: substitution

\documentclass[../../../include/open-logic-section]{subfiles}

\begin{document}

\olfileid{nml}{syn}{sub}

\olsection{ఏకకాల ప్రతిస్థాపన}

\tetoken{సూత్రం}{!!a{formula}}~$!A$కు ఒక
\emph{నిదర్శనం} అంటే, $!A$లోని ఒక
\tetoken{ప్రతిజ్ఞావాక్య చరం}{!!a{propositional variable}}
యొక్క అన్ని ఘటనలను మరొక
\tetoken{సూత్రంతో}{!!{formula}} ప్రతిస్థాపించి
వచ్చిన ఫలితం. చెల్లుబాటును,
\tetoken{నిరూప్యతను}{!!{derivability}}
చర్చించేటప్పుడు \tetoken{సూత్రాల}{!!{formula}s}
నిదర్శనాలను తరచుగా ప్రస్తావిస్తాం.
కాబట్టి ఈ భావాన్ని ఖచ్చితంగా
నిర్వచించడం ఉపయోగకరం.

\begin{defn}\ollabel{def:subst-inst}
  మోడల్ \tetoken{సూత్రం}{!!{formula}}~$!A$లోని
  \tetoken{ప్రతిజ్ఞావాక్య చరాలన్నీ}{!!{propositional
    variable}s} $p_1$, \dots, $p_n$లలోనే
  ఉండి, $!D_1$, \dots, $!D_n$లు కూడా
  మోడల్ \tetoken{సూత్రాలైతే}{!!{formula}s},
  $\SSubst{!A}{\subst{!D_1}{p_1}, \dots, \subst{!D_n}{p_n}}$ను
  $!A$లో ప్రతి $p_i$ స్థానంలో $!D_i$ని
  ఏకకాలంలో ప్రతిస్థాపించి వచ్చిన ఫలితంగా
  నిర్వచిస్తాం. అధికారికంగా, $!A$పై
  ఆగమనంతో నిర్వచనం ఇలా ఉంటుంది:
  \begin{enumerate}
    \tagitem{prvFalse}{\indcase{!A}{\lfalse}{%
        $\SSubst{\indfrm}{\subst{!D_1}{p_1}, \dots,
          \subst{!D_n}{p_n}}$ అనేది $\lfalse$.}}{}
    \tagitem{prvTrue}{\indcase{!A}{\ltrue}{%
        $\SSubst{\indfrm}{\subst{!D_1}{p_1}, \dots,
          \subst{!D_n}{p_n}}$ అనేది $\ltrue$.}}{}
    \item \indcase{!A}{q}{$i
      = 1$, \dots,~$n$ అన్నింటికీ
      $q \not\ident p_i$ అయితే,
      $\SSubst{\indfrm}{\subst{!D_1}{p_1}, \dots,
        \subst{!D_n}{p_n}}$ అనేది $q$.}
    \item \indcase{!A}{p_i}{$\SSubst{\indfrm}{\subst{!D_1}{p_1},
        \dots, \subst{!D_n}{p_n}}$ అనేది $!D_i$.}
    \tagitem{prvNot}{\indcase{!A}{\lnot !B}{%
        $\SSubst{\indfrm}{\subst{!D_1}{p_1}, \dots, \subst{!D_n}{p_n}}$
        అనేది
        $\lnot \SSubst{!B}{\subst{!D_1}{p_1}, \dots, \subst{!D_n}{p_n}}$.}}{}
    \tagitem{prvAnd}{\indcase{!A}{(!B \land
        !C)}{$\SSubst{\indfrm}{\subst{!D_1}{p_1}, \dots,
          \subst{!D_n}{p_n}}$ అనేది
        \[(\SSubst{!B}{\subst{!D_1}{p_1}, \dots, \subst{!D_n}{p_n}} \land
          \SSubst{!C}{\subst{!D_1}{p_1}, \dots, \subst{!D_n}{p_n}}).\]}}{}
    \tagitem{prvOr}{\indcase{!A}{(!B \lor
          !C)}{$\SSubst{\indfrm}{\subst{!D_1}{p_1}, \dots,
          \subst{!D_n}{p_n}}$ అనేది
        \[(\SSubst{!B}{\subst{!D_1}{p_1}, \dots, \subst{!D_n}{p_n}} \lor
          \SSubst{!C}{\subst{!D_1}{p_1}, \dots, \subst{!D_n}{p_n}}).\]}}{}
    \tagitem{prvIf}{\indcase{!A}{(!B \lif
          !C)}{$\SSubst{\indfrm}{\subst{!D_1}{p_1}, \dots,
          \subst{!D_n}{p_n}}$ అనేది
        \[(\SSubst{!B}{\subst{!D_1}{p_1}, \dots, \subst{!D_n}{p_n}} \lif
          \SSubst{!C}{\subst{!D_1}{p_1}, \dots, \subst{!D_n}{p_n}}).\]}}{}
    \tagitem{prvIff}{\indcase{!A}{(!B \liff
          !C)}{$\SSubst{\indfrm}{\subst{!D_1}{p_1}, \dots,
          \subst{!D_n}{p_n}}$ అనేది
        \[(\SSubst{!B}{\subst{!D_1}{p_1}, \dots, \subst{!D_n}{p_n}} \liff
          \SSubst{!C}{\subst{!D_1}{p_1}, \dots, \subst{!D_n}{p_n}}).\]}
      \sourcecorrection{OLTENMLSYN-001}{ఆంగ్ల మూలంలో
        ద్విసోపాధిక సందర్భానికి prvIf గుర్తు
        ఉంది; పక్కనున్న \liff సూత్రానుసారం
        prvIffగా సరిచేశాం.}}{}
    \tagitem{prvBox}{\indcase{!A}{\Box !B}{%
        $\SSubst{\indfrm}{\subst{!D_1}{p_1}, \dots, \subst{!D_n}{p_n}}$
        అనేది $\Box
        \SSubst{!B}{\subst{!D_1}{p_1}, \dots, \subst{!D_n}{p_n}}.$}
      \sourcecorrection{OLTENMLSYN-002}{మూలంలో
        ఈ బాక్స్ సందర్భం సాధారణ జాబితా
        అంశంగా ఉంది; భాషా నిర్వచనంలోని
        prvBox షరతుకు అనుగుణంగా
        అదే సందర్భాన్ని ఆ గుర్తుతో
        కట్టాం.}}{}
    \tagitem{prvDiamond}{\indcase{!A}{\Diamond !B}{%
        $\SSubst{\indfrm}{\subst{!D_1}{p_1}, \dots, \subst{!D_n}{p_n}}$
        అనేది $\Diamond
        \SSubst{!B}{\subst{!D_1}{p_1}, \dots, \subst{!D_n}{p_n}}.$}}{}
  \end{enumerate}
  \tetoken{సూత్రం}{!!{formula}}
  $\SSubst{!A}{\subst{!D_1}{p_1}, \dots,
    \subst{!D_n}{p_n}}$ను $!A$కి
  \emph{ప్రతిస్థాపన నిదర్శనం} అంటాం.
\end{defn}

\begin{ex}
  $!A$ అనేది $p_1 \lif \Box(p_1 \land p_2)$,
  $!D_1$ అనేది $\Diamond(p_2 \lif p_3)$,
  $!D_2$ అనేది $\lnot\Box p_1$
  అనుకుందాం. అప్పుడు
  $\SSubst{!A}{\subst{!D_1}{p_1}, \subst{!D_2}{p_2}}$
  ఇలా ఉంటుంది:
  \begin{align*}
    \Diamond(p_2 \lif p_3) &
    \lif \Box(\Diamond(p_2 \lif p_3) \land \lnot\Box p_1)
    \intertext{అయితే
      $\SSubst{!A}{\subst{!D_2}{p_1}, \subst{!D_1}{p_2}}$
      ఇలా ఉంటుంది:}
    \lnot\Box p_1 &
    \lif \Box(\lnot\Box p_1 \land \Diamond(p_2 \lif p_3))
    \intertext{ఏకకాల ప్రతిస్థాపన సాధారణంగా
      వరుస ప్రతిస్థాపనతో సమానం కాదని గమనించండి.
      ఉదాహరణకు, పైన ఉన్న
      $\SSubst{!A}{\subst{!D_1}{p_1}, \subst{!D_2}{p_2}}$ను
      $\Subst{(\Subst{!A}{!D_1}{p_1})}{!D_2}{p_2}$తో
      పోల్చండి; రెండవది ఇలా ఉంటుంది:}
    \Diamond(p_2 \lif p_3) & \Subst{\lif \Box(\Diamond(p_2 \lif p_3) \land p_2)}{\lnot\Box p_1}{p_2}, \text{ అంటే,}\\
    \Diamond(\lnot\Box p_1 \lif p_3) &
    \lif \Box(\Diamond(\lnot\Box p_1
    \lif p_3) \land \lnot\Box p_1)\\
    \intertext{అలాగే
      $\Subst{(\Subst{!A}{!D_2}{p_2})}{!D_1}{p_1}$తో
      పోలిస్తే:}
    p_1 & \lif \Subst{\Box(p_1 \land \lnot\Box p_1)}{\Diamond(p_2 \lif p_3)}{p_1}, \text{ అంటే,}\\
    \Diamond(p_2 \lif p_3) &
    \lif \Box(\Diamond(p_2
    \lif p_3) \land \lnot\Box\Diamond(p_2 \lif p_3)).
  \end{align*}
\end{ex}

\end{document}
